\documentclass[11pt,a4paper]{article}

\usepackage[utf8]{inputenc}
\usepackage[T1]{fontenc}
\usepackage{amsmath,amssymb,amsthm}
\usepackage{algorithm}
\usepackage{algorithmic}
\usepackage{graphicx}
\usepackage{hyperref}
\usepackage{booktabs}
\usepackage{listings}
\usepackage[margin=1in]{geometry}

\title{\textbf{CUDA-zkML: GPU-Accelerated Zero-Knowledge Proofs\\
for Neural Network Inference}}

\author{George David Tsitlauri\\
\textit{Dept. of Informatics \& Telecommunications}\\
\textit{University of Thessaly, Greece}\\
gdtsitlauri@gmail.com}

\date{2026}

\begin{document}

\maketitle

% ============================================================
\begin{abstract}
We present CUDA-zkML, a GPU-accelerated system for generating
zero-knowledge proofs of neural network inference. Our system proves
that a neural network was correctly executed on given inputs without
revealing the model weights or private inputs. Built as a CUDA-native
implementation targeting NVIDIA consumer GPUs (GTX 1650, SM 7.5),
CUDA-zkML achieves significant speedups over CPU-based zkML tools
by parallelizing the computationally intensive multi-scalar multiplication
(MSM) and number-theoretic transform (NTT) operations that dominate
Groth16 proof generation. We implement complete BN254 finite field
arithmetic in CUDA, a quantized neural network inference engine over
finite fields, and a full Groth16 proving system. On an NVIDIA GTX 1650
(4GB VRAM), the end-to-end demo path measures 53.66 ms inference,
1381.30 ms setup, 3982.29 ms proving, and 225.18 ms native verification
for an MNIST MLP, while also supporting local EVM verification through a
Solidity Groth16 verifier harness. These results demonstrate the
feasibility of GPU-accelerated zkML on commodity hardware with just 4GB
of VRAM.
\end{abstract}

% ============================================================
\section{Introduction}

Zero-knowledge machine learning (zkML) enables a prover to demonstrate
that a neural network was correctly executed without revealing sensitive
information such as model weights (intellectual property protection)
or input data (privacy preservation). This has applications in:

\begin{itemize}
  \item \textbf{Model verification}: Proving that a specific model
    produced a given prediction
  \item \textbf{Privacy-preserving inference}: Running ML inference
    on private data with public verifiability
  \item \textbf{On-chain ML}: Verifying neural network outputs in
    smart contracts without executing the model on-chain
\end{itemize}

The primary bottleneck in zkML is proof generation time. A Groth16
proof for even a small neural network requires millions of elliptic
curve operations. CPU-based systems like EZKL~\cite{ezkl} require
seconds to minutes for small models. GPU acceleration can dramatically
reduce this time by parallelizing the dominant operations:
multi-scalar multiplication (MSM) and number-theoretic transform (NTT).

Our contributions:
\begin{enumerate}
  \item A complete CUDA implementation of BN254 finite field and
    elliptic curve arithmetic optimized for the Turing architecture
  \item GPU-parallel Pippenger MSM achieving significant speedup
    over sequential CPU implementations
  \item A quantized neural network inference engine that operates
    entirely in finite field arithmetic, with architecture sidecars for
    ONNX-derived feedforward models, CIFAR-10-sized flattened MLPs,
    and transformer-style attention blocks including stacked multi-head variants
  \item A Groth16 proving system with GPU-accelerated MSM and NTT
  \item A local Solidity/EVM verification workflow exercised against the
    generated proof artifacts
  \item Demonstration of zkML feasibility on a consumer GPU
    (NVIDIA GTX 1650, 4GB VRAM)
\end{enumerate}

% ============================================================
\section{Background}

\subsection{Zero-Knowledge Proofs}

A zero-knowledge proof system allows a prover to convince a verifier
that a statement is true without revealing any additional information.
In the context of zkML, the statement is ``I correctly ran neural
network $f$ on input $x$ and obtained output $y$''.

\textbf{Groth16}~\cite{groth16} is a pairing-based zk-SNARK
(Succinct Non-interactive Argument of Knowledge) with:
\begin{itemize}
  \item Constant-size proofs: 3 group elements ($\sim$256 bytes)
  \item Constant verification time: 3 pairings
  \item Trusted setup required (per-circuit)
  \item Proof generation: $O(n \log n)$ field operations
\end{itemize}

\subsection{zkML}

zkML converts neural network computation into an arithmetic circuit,
generates a witness (all intermediate values), and produces a ZK proof.
The key challenge is the circuit size: a matrix multiplication of
dimension $M \times N$ produces $O(MN)$ multiplication constraints.

\subsection{GPU Parallelism for ZK Proofs}

The two most computationally expensive operations in Groth16 proving are:

\textbf{Multi-Scalar Multiplication (MSM):}
$\text{MSM}(\{s_i\}, \{P_i\}) = \sum_{i=1}^n s_i \cdot P_i$,
where $s_i$ are scalars and $P_i$ are elliptic curve points.
This is embarrassingly parallel using Pippenger's bucket method.

\textbf{Number-Theoretic Transform (NTT):}
The finite-field analogue of FFT, used for polynomial multiplication
in $O(n \log n)$. NTT has a natural butterfly-network structure
suitable for GPU parallel execution.

% ============================================================
\section{System Design}

\subsection{Architecture}

\begin{verbatim}
+------------------+     +------------------+
| Model (ONNX/PT)  |---->| Quantizer        |
+------------------+     | float -> Fp      |
                          +--------+---------+
                                   |
+------------------+     +---------v---------+
| Input Data       |---->| Field Inference   |
+------------------+     | (GPU, in Fp)      |
                          +--------+---------+
                                   |
                          +--------v---------+
                          | Witness Generator |
                          +--------+---------+
                                   |
                          +--------v---------+
                          | Groth16 Prover   |
                          | (GPU: MSM + NTT) |
                          +--------+---------+
                                   |
                          +--------v---------+
                          | Proof (A, B, C)  |
                          +------------------+
\end{verbatim}

\subsection{Memory Management}

Targeting the GTX 1650 with 4GB VRAM, we enforce:
\begin{itemize}
  \item Maximum single allocation: 1GB
  \item Usable VRAM budget: 3.5GB (500MB headroom)
  \item Streaming MSM for $> 1\text{M}$ points
  \item NTT maximum size: $2^{24}$ elements with ping-pong buffers
  \item \texttt{cudaMemGetInfo()} checks before large allocations
\end{itemize}

% ============================================================
\section{Implementation}

\subsection{BN254 Field Arithmetic}

We implement 254-bit modular arithmetic using Montgomery multiplication
in CUDA. Each field element is stored as 4 $\times$ \texttt{uint64\_t}
in Montgomery form ($aR \bmod p$ where $R = 2^{256}$).

Montgomery multiplication uses the CIOS (Coarsely Integrated Operand
Scanning) algorithm, requiring $4 \times 4 = 16$ 64-bit multiplications
per field multiplication.

\subsection{Pippenger MSM Kernel}

\begin{algorithm}
\caption{GPU-Parallel Pippenger MSM}
\begin{algorithmic}[1]
\REQUIRE Points $\{P_i\}_{i=1}^n$, scalars $\{s_i\}_{i=1}^n$, window size $c$
\ENSURE $R = \sum_{i=1}^n s_i \cdot P_i$
\STATE $W \leftarrow \lceil 254/c \rceil$ \COMMENT{number of windows}
\STATE $K \leftarrow 2^c - 1$ \COMMENT{buckets per window}
\STATE Initialize buckets $B[w][k] \leftarrow \mathcal{O}$ for $w \in [W], k \in [K]$
\STATE \textbf{// Phase 1: Bucket accumulation (parallel over points)}
\FOR{each point $i$ \textbf{in parallel}}
  \FOR{each window $w = 0, \ldots, W-1$}
    \STATE $b \leftarrow (s_i \gg (w \cdot c)) \bmod 2^c$
    \IF{$b \neq 0$}
      \STATE $B[w][b-1] \leftarrow B[w][b-1] + P_i$
    \ENDIF
  \ENDFOR
\ENDFOR
\STATE \textbf{// Phase 2: Bucket reduction (parallel over windows)}
\FOR{each window $w$ \textbf{in parallel}}
  \STATE $\text{running} \leftarrow \mathcal{O}$, $\text{sum}_w \leftarrow \mathcal{O}$
  \FOR{$k = K$ \textbf{down to} $1$}
    \STATE $\text{running} \leftarrow \text{running} + B[w][k-1]$
    \STATE $\text{sum}_w \leftarrow \text{sum}_w + \text{running}$
  \ENDFOR
\ENDFOR
\STATE \textbf{// Phase 3: Window combination (sequential)}
\STATE $R \leftarrow \mathcal{O}$
\FOR{$w = W-1$ \textbf{down to} $0$}
  \STATE $R \leftarrow 2^c \cdot R$ \COMMENT{$c$ doublings}
  \STATE $R \leftarrow R + \text{sum}_w$
\ENDFOR
\RETURN $R$
\end{algorithmic}
\end{algorithm}

The optimal window size is $c \approx \log_2 n$ for $n$ points.
For $n = 2^{16}$, $c = 16$ gives $2^{16} - 1 = 65535$ buckets per window,
with $\lceil 254/16 \rceil = 16$ windows.

\subsection{Quantized Neural Network Inference}

We use symmetric quantization:
\[
  x_q = \text{round}(x \cdot s), \quad
  s = \frac{2^{b-1} - 1}{\max |x|}
\]

Integer values are mapped to field elements:
$x_q \mapsto x_q \bmod p$ (negative values use additive inverse).

Matrix multiplication in the field is exact (no rounding):
\[
  y_i = \sum_j W_{ij} \cdot x_j \pmod{p}
\]

ReLU is approximated by a degree-3 polynomial to maintain
arithmetic circuit compatibility.

The current inference stack supports MNIST-sized MLPs, CIFAR-10 flattened
MLPs, a tiny transformer block family (4 tokens $\times$ 8 hidden
channels) including stacked multi-head attention via architecture sidecars,
and ONNX models that lower to linear, ReLU, final softmax, and
the repository's self-attention layer description. Model conversion emits an
architecture sidecar that allows the prover CLI to reconstruct the network
topology without a hardcoded layer schedule.

\subsection{Witness Generation}

The witness captures every intermediate value during inference:
$w = [1, \text{outputs}, \text{inputs}, \text{weights}, \text{intermediates}]$.

For an MLP with layers of sizes $[784, 128, 10]$:
\begin{itemize}
  \item Public: 10 output values
  \item Private: 784 input pixels + $784 \times 128 + 128 \times 10 = 101,\!632$ weights
  \item Intermediates: partial products and activations
\end{itemize}

% ============================================================
\section{Evaluation}

\subsection{Hardware}

\begin{tabular}{ll}
\toprule
Component & Specification \\
\midrule
GPU & NVIDIA GTX 1650 (Turing, SM 7.5) \\
CUDA Cores & 896 \\
VRAM & 4 GB GDDR6, 112 GB/s bandwidth \\
CPU & AMD Ryzen 5 5600H \\
RAM & 16 GB DDR4 \\
CUDA Version & 12.x \\
\bottomrule
\end{tabular}

\subsection{Benchmark Results}

\begin{tabular}{lcccc}
\toprule
System & Status & Prove (ms) & Verify (ms) & Proof Size \\
\midrule
CUDA-zkML (GTX 1650) & measured & 3982.29 & 225.18 & 256 B \\
CUDA-zkML Tiny Transformer & measured & 207.14 & 263.35 & 256 B \\
EZKL (WSL Ubuntu live run) & measured & 1812.74 & 13.26 & 18209 B \\
Orion (WSL filtered relu\_i8) & measured & 456930.00 & N/A & N/A \\
\bottomrule
\end{tabular}

Additional measured timings for CUDA-zkML on the target laptop are
53.66 ms inference and 1381.30 ms trusted setup for the MNIST demo path.
The tiny transformer path measures 44.84 ms inference and 174.79 ms setup,
demonstrating that the same proving pipeline now covers a compact attention
architecture in addition to feedforward models. The local CPU baselines
included in the repository are inference-only references
(0.09 ms float32, 1.09 ms quantized-field) and therefore are not directly
comparable to proving systems. EZKL was measured successfully in a WSL Ubuntu
environment on the same laptop, yielding 5.07 ms witness generation,
1547.85 ms setup, 1812.74 ms proving, and 13.26 ms verification for the
repository's tiny reference ONNX model. The native Windows EZKL wheel still
fails inside \texttt{setup(...)} with a runtime \texttt{NotPresent} panic, so
WSL/Linux remains the reliable path on this machine. Orion was measured from a
passing filtered archived test target (\texttt{relu\_i8}) in a local WSL Ubuntu
checkout. The full archived Orion workspace test suite still reports upstream
failures on this setup, so the filtered row is the stable comparison point
captured in the repository artifacts.

\subsection{On-Chain Verification}

We also exercise the Solidity verifier in a local EVM environment using
\texttt{eth-tester} and \texttt{py-evm}. For the generated demo proof,
deployment succeeds, the verification key is uploaded on-chain, and
\texttt{verifyProofView(...)} returns \texttt{true}. The measured local
gas costs are 2,115,563 gas for contract deployment and 894,575 gas for
verification-key upload. This workflow provides an executable bridge from
the native prover artifacts to EVM-compatible verification.

% ============================================================
\section{Related Work}

\textbf{EZKL}~\cite{ezkl}: A Rust/Python framework for generating
ZK proofs of ONNX model inference. Operates on CPU with
Halo2 proving system.

\textbf{Orion} (Giza): Cairo-based zkML framework targeting
StarkNet's STARK-based proofs.

\textbf{Daniel Kang et al.}: ``Scaling up Trustless DNN Inference
with Zero-Knowledge Proofs'' demonstrated zkML feasibility but
with high proving costs.

\textbf{GPU ZK Libraries}: icicle (Ingonyama), Sppark (Supranational),
and Bellman-CUDA provide GPU-accelerated primitives for ZK proofs,
but none provide an end-to-end zkML pipeline.

Our work differs by providing a complete, self-contained CUDA
implementation from field arithmetic to proof generation, specifically
optimized for consumer GPUs, together with an executable native-to-EVM
artifact flow.

% ============================================================
\section{Conclusion}

We presented CUDA-zkML, a GPU-accelerated zero-knowledge proof
system for neural network inference. By implementing the entire
proving pipeline in CUDA—from finite field arithmetic through
Pippenger MSM to Groth16 proof generation—we demonstrate that
zkML is feasible on consumer GPUs with limited VRAM.

Future work includes:
\begin{itemize}
  \item Support for convolutional and full transformer attention architectures
  \item PLONK proving system for universal (no per-circuit) setup
  \item Multi-GPU support for larger models
  \item Integration with ONNX Runtime for seamless model conversion
  \item Formal verification of the CUDA arithmetic kernels
\end{itemize}

% ============================================================
\begin{thebibliography}{9}

\bibitem{groth16}
J.~Groth, ``On the size of pairing-based non-interactive arguments,''
in \textit{EUROCRYPT 2016}, pp.~305--326, Springer, 2016.

\bibitem{ezkl}
EZKL, ``Easy Zero-Knowledge inference,''
\url{https://github.com/zkonduit/ezkl}, 2023.

\bibitem{pippenger}
B.~Pippenger, ``On the evaluation of powers and monomials,''
\textit{SIAM J. Comput.}, vol.~9, no.~2, pp.~230--250, 1980.

\bibitem{bn254}
P.~Barreto and M.~Naehrig, ``Pairing-friendly elliptic curves
of prime order,'' in \textit{SAC 2005}, pp.~319--331, 2005.

\bibitem{montgomery}
P.~Montgomery, ``Modular multiplication without trial division,''
\textit{Mathematics of Computation}, vol.~44, pp.~519--521, 1985.

\bibitem{bellman}
Zcash, ``bellman: zk-SNARK library,''
\url{https://github.com/zkcrypto/bellman}, 2017.

\bibitem{arkworks}
arkworks contributors, ``arkworks: An ecosystem for developing
and programming with zkSNARKs,''
\url{https://github.com/arkworks-rs}, 2021.

\bibitem{icicle}
Ingonyama, ``ICICLE: GPU library for zero-knowledge,''
\url{https://github.com/ingonyama-zk/icicle}, 2023.

\bibitem{kang2022}
D.~Kang, T.~Hashimoto, I.~Stoica, and Y.~Sun,
``Scaling up trustless DNN inference with zero-knowledge proofs,''
\textit{arXiv:2210.08674}, 2022.

\end{thebibliography}

\end{document}
